% Continuation of the frozen domain-theory-and-recursive-models packet.
\section{Canonical recursion, invariance and relational reasoning}
\label{sec:p10-domain-recursion-invariance}

An inverse pair solving a domain equation does not by itself establish
minimality, initiality or finality. These conclusions require the
canonical construction and its least-endomorphism property. This
continuation develops that dependency from the corrected and expanded
Abramsky--Jung author presentation, including its separate strictness
regimes \cite[Sections~5.2.3--5.4, pp.~73--84]
{AbramskyJungCorrectedExpandedDomainTheory}.
The underlying directed, function-space and bilimit constructions are
Theorems~\ref{thm:p10-domain-function-space},
\ref{thm:p10-domain-least-fixpoint},\ref{thm:p10-domain-bilimit}
and~\ref{thm:p10-domain-recursive-solution}.

\begin{hbconvention}{def:p10-domain-invariance-category}{Objects, hom dcpo's and the canonical seed}
The default category \(\mathcal D\) is
\(\mathsf{DCPO}_{\bot}\), with set-sized pointed dcpo's and all
Scott-continuous maps under the previously declared universe convention.
One may instead use a full subcategory containing the one-point
\(I\) and closed under the actual expanding-sequence bilimits;
its hom objects still contain all continuous maps between its objects.
Write \(\mathcal D_{\bot!}\) for the same objects with strict maps.
Its hom dcpo's are the strict function spaces, not an assumption that
an unrestricted functor preserves strict maps.
All functors below are locally continuous on the displayed hom dcpo's.
For a mixed functor this includes correctly typed joint continuity in
the contravariant and covariant hom variables. Products of finitely many
such categories use coordinatewise hom orders and bilimits.
Every \(\operatorname{FIX}\) in this section starts at \(I\),
not at an arbitrary seed from Theorem~\ref{thm:p10-domain-recursive-solution}.
All colimits used are countable expanding-sequence bilimits; no
unasserted arbitrary colimit, universe enlargement or effective
presentation is required \cite[pp.~70--74]
{AbramskyJungCorrectedExpandedDomainTheory}.
\end{hbconvention}

\subsection{Least comparisons and covariant universal properties}

\begin{hbdefinition}{def:p10-domain-invariance-algebra}{Algebras, coalgebras and solution comparisons}
For a covariant \(F\), an \(F\)-algebra has structure
\(a:F(A)\to A\); an algebra homomorphism \(h:A\to B\)
satisfies \(ha=bF(h)\). A coalgebra has structure
\(c:A\to F(A)\); a coalgebra homomorphism into
\(d:B\to F(B)\) satisfies \(dh=F(h)c\).
A solution \((A,i)\) has an actual isomorphism \(i:F(A)\to A\).
An isomorphism of solutions includes preservation of these structure
maps; it is not merely an untyped bijection of carriers.
For the canonical solution \(D\), write
\(i_D=\operatorname{fold}\), \(i_D^{-1}=\operatorname{unfold}\),
and \(q_n=e_np_n\) for its bilimit approximating idempotents.
\cite[Definition~5.1.1; pp.~69--70,74--77]
{AbramskyJungCorrectedExpandedDomainTheory}.
\end{hbdefinition}

\begin{hbtheorem}{thm:p10-domain-invariance-least-hom}{proved=annex:p10.domain-invariance.least-hom}{Least homomorphisms and their composition}
Let \(F:\mathcal D\to\mathcal D\) be locally continuous,
\(i:F(A)\to A\) an isomorphism, and \(a:F(C)\to C\)
any continuous algebra structure. Define
\[
 \phi_{C,A}(u)=aF(u)i^{-1},\qquad
 h_{C,A}=\operatorname{fix}(\phi_{C,A})
                       \quad\text{in }[A\longrightarrow C].
\]
This is the least algebra homomorphism. For another solution \((B,j)\),
\(h_{C,A}=h_{C,B}h_{B,A}\). Consequently
\(\mu_A=h_{A,A}\) is an idempotent below identity, and is strict.
\cite[Lemma~5.3.1, pp.~74--75]
{AbramskyJungCorrectedExpandedDomainTheory}.
\end{hbtheorem}

\begin{hbtheorem}{thm:p10-domain-invariance-covariant}{proved=annex:p10.domain-invariance.covariant}{Canonical identity, invariance and minimality}
For the canonical \(D=\operatorname{FIX}(F)\),
\[
 \theta_D(u)=i_DF(u)i_D^{-1},\qquad
 \theta_D^n(c_{\bot})=q_n,\qquad
 \operatorname{fix}(\theta_D)=1_D.
\]
For any solution \((A,i)\), its least comparisons
\(e=h_{A,D}:D\to A\), \(p=h_{D,A}:A\to D\)
satisfy \(pe=1_D\) and \(ep=\mu_A\sqsubseteq1_A\).
Thus the canonical solution is a subdomain of every solution.
The following are equivalent: \((A,i)\) is isomorphic as a solution
to the canonical solution; \(\mu_A=1_A\); and identity is the
only strict algebra endomorphism of \((A,i)\).
\cite[Theorems~5.3.2--3, pp.~75--76]
{AbramskyJungCorrectedExpandedDomainTheory}.
\end{hbtheorem}

\begin{hbtheorem}{thm:p10-domain-invariance-initial-final}{proved=annex:p10.domain-invariance.initial-final}{Initiality and finality in their different regimes}
If \(F:\mathcal D_{\bot!}\to\mathcal D_{\bot!}\) is
locally continuous, canonical fold is an initial \(F\)-algebra
in the strict-map category. Its unique strict homomorphism to a
strict algebra \(a:F(A)\to A\) is the least fixed point of
\(u\mapsto aF(u)\operatorname{unfold}\) on the strict hom dcpo.
If instead \(F:\mathcal D\to\mathcal D\) is locally continuous
on unrestricted continuous maps, canonical unfold is a final
coalgebra for \emph{all} continuous coalgebra structures
\(c:A\to F(A)\), including non-strict ones. Its unique
homomorphism is
\[
 \operatorname{fix}\bigl(u\mapsto\operatorname{fold}F(u)c\bigr)
                              \quad\text{in }[A\longrightarrow D].
\]
The second assertion does not extend a functor known only on strict
maps to non-strict maps. In the strict-only regime it gives finality
for strict coalgebras and strict homomorphisms.
\cite[Theorems~5.3.4--5; Exercise~5.4.11(6), pp.~77,83]
{AbramskyJungCorrectedExpandedDomainTheory}.
\end{hbtheorem}

\subsection{Mixed variance and the free comparison property}

\begin{hbdefinition}{def:p10-domain-invariance-mixed-comparison}{Paired comparisons for mixed solutions}
For \(F:\mathcal D^{\mathrm{op}}\times\mathcal D\to\mathcal D\),
solutions have \(i:F(A,A)\to A\), \(j:F(B,B)\to B\).
A comparison pair \(h:A\to B,k:B\to A\) satisfies
\[
 hi=jF(k,h),\qquad kj=iF(h,k).
\]
These are two typed equations, not an ordinary covariant
\(F\)-algebra homomorphism for the diagonal object assignment.
A mixed endomorphism of \((A,i)\) satisfies \(hi=iF(h,h)\).
\cite[Lemma~5.3.6, pp.~77--78]
{AbramskyJungCorrectedExpandedDomainTheory}.
\end{hbdefinition}

\begin{hbtheorem}{thm:p10-domain-invariance-least-pairs}{proved=annex:p10.domain-invariance.least-pairs}{Least pairs, exact composition and the mixed kernel}
The least comparison pair is the least fixed point of
\[
 \Phi_{B,A}(u,v)
       =(jF(v,u)i^{-1},\ iF(u,v)j^{-1})
\]
on \([A\longrightarrow B]\times[B\longrightarrow A]\).
For a third solution, composing forward maps forward and backward
maps backward gives the least comparison pair between the outside
solutions. For \(A=B\) the two least maps agree and equal
\[
 \mu_A=\operatorname{fix}\bigl(u\mapsto iF(u,u)i^{-1}\bigr).
\]
This mixed kernel is idempotent, strict and below identity.
\cite[Lemma~5.3.6 and its corrected-author comment, p.~78]
{AbramskyJungCorrectedExpandedDomainTheory}.
\end{hbtheorem}

\begin{hbtheorem}{thm:p10-domain-invariance-mixed}{proved=annex:p10.domain-invariance.mixed}{Mixed canonical identity and minimality}
Let \(D=\operatorname{FIX}(\widehat F\circ\Delta)\), using the
actual embedding functor of Theorem~\ref{thm:p10-domain-mixed-embedding}.
Then for \(\theta_D(u)=\operatorname{fold}F(u,u)\operatorname{unfold}\),
\(\theta_D^n(c_{\bot})=q_n\) and \(\mu_D=1_D\).
Its least comparison pair with any mixed solution \((A,i)\)
is an embedding--projection pair with composite on \(A\) equal
to \(\mu_A\). Moreover canonical-solution isomorphism,
\(\mu_A=1_A\), and uniqueness of the strict mixed endomorphism
are equivalent. In the uniqueness proof, \(u\mapsto huh\)
first gives \(h^2=1\); the least-identity order comparison is
also used to conclude \(h=1\).
\cite[Theorems~5.3.7--8, pp.~78--79]
{AbramskyJungCorrectedExpandedDomainTheory}.
\end{hbtheorem}

\begin{hbtheorem}{thm:p10-domain-invariance-free-mixed}{proved=annex:p10.domain-invariance.free-mixed}{The actual free mixed-variant universal property}
Suppose \(F:\mathcal D_{\bot!}^{\mathrm{op}}\times
\mathcal D_{\bot!}\to\mathcal D_{\bot!}\) is locally
continuous and \(D\) its canonical mixed solution. For strict
\(f:A\to F(B,A)\), \(g:F(A,B)\to B\), there is a
\emph{unique strict} pair \(h:A\to D,k:D\to B\) with
\[
 \operatorname{unfold}h=F(k,h)f,\qquad
 k\operatorname{fold}=gF(h,k).
\]
It is the least fixed point on the strict product hom dcpo of
\[
 (u,v)\longmapsto
 (\operatorname{fold}F(v,u)f,\ gF(u,v)\operatorname{unfold}).
\]
Uniqueness is proved by the canonical approximants \(q_n\),
not by the mere invertibility of fold and unfold.
\cite[Theorem~5.3.9, p.~79]{AbramskyJungCorrectedExpandedDomainTheory}.
\end{hbtheorem}

\subsection{Parameters, simultaneous equations and resumptions}

\begin{hbtheorem}{thm:p10-domain-invariance-parameters}{proved=annex:p10.domain-invariance.parameters}{The parameter solution is an actual locally continuous functor}
Let \(F:\mathcal C_{\bot!}\times\mathcal D_{\bot!}\to
\mathcal D_{\bot!}\) be locally continuous. Put
\(F_A(X)=F(A,X)\), \(M(A)=\operatorname{FIX}(F_A)\), and
\(X_n(A)=F_A^n(I)\). For a strict \(u:A\to A'\), define
\[
 u_0=1_I,\quad u_{n+1}=F(u,u_n),\qquad
 M(u)=\bigsqcup_n^{\uparrow}e_n^{A'}u_np_n^A.
\]
These maps are strict, the supremum is directed, and this gives
a locally continuous functor. Its actual structure isomorphisms
\(\alpha_A:F(A,M(A))\to M(A)\) satisfy
\[
 M(u)\alpha_A=\alpha_{A'}F(u,M(u)),\qquad
 \alpha_{A'}^{-1}M(u)=F(u,M(u))\alpha_A^{-1}.
\]
\cite[Proposition~5.2.7, pp.~73--74]
{AbramskyJungCorrectedExpandedDomainTheory}.
\end{hbtheorem}

\begin{hbtheorem}{thm:p10-domain-invariance-mixed-parameters}{proved=annex:p10.domain-invariance.mixed-parameters}{Explicitly paired mixed parameters}
Let \(F:\mathcal C_{\bot!}^{\mathrm{op}}\times
\mathcal C_{\bot!}\times\mathcal D_{\bot!}^{\mathrm{op}}
\times\mathcal D_{\bot!}\to\mathcal D_{\bot!}\) be locally
continuous. For parameter objects \(P,Q\), start \(X_0=Y_0=I\)
and iterate
\[
 X_{n+1}=F(P,Q,Y_n,X_n),\quad
 Y_{n+1}=F(Q,P,X_n,Y_n).
\]
The mixed embedding action gives actual expanding sequences.
Their bilimits \(M(P,Q)\), \(M(Q,P)\) satisfy
\[
 M(P,Q)\cong F(P,Q,M(Q,P),M(P,Q)).
\]
For \(u:P'\to P,v:Q\to Q'\), define
\[
 r_0=s_0=1_I,\quad
 r_{n+1}=F(u,v,s_n,r_n),\quad
 s_{n+1}=F(v,u,r_n,s_n).
\]
The first sequence induces \(M(u,v):M(P,Q)\to M(P',Q')\);
the second induces the reversed map \(M(v,u)\).
The induced maps use the same bilimit supremum formula as the
preceding theorem and define a locally continuous mixed functor.
The structure maps are natural with the typed paired action
\(F(u,v,M(v,u),M(u,v))\). This makes the negative recursive
slot explicit instead of substituting a covariant action there.
This is the four-slot extension of the source's parameter
construction and its prescribed variance handling
\cite[pp.~72--74; Exercise~5.4.11(9), p.~84]
{AbramskyJungCorrectedExpandedDomainTheory}.
\end{hbtheorem}

\begin{hbtheorem}{thm:p10-domain-invariance-bekic}{proved=annex:p10.domain-invariance.bekic}{Beki\v c elimination, for values and for canonical domains}
For pointed dcpo's \(D,E\) and continuous
\(f:D\times E\to D,g:D\times E\to E\), set
\[
 h(y)=\operatorname{fix}(x\mapsto f(x,y)),\quad
 k(y)=g(h(y),y),\quad d=\operatorname{fix}(k),\quad c=h(d).
\]
Then \((c,d)=\operatorname{fix}\langle f,g\rangle\).
For locally continuous strict functors
\(P:\mathcal C_{\bot!}\times\mathcal E_{\bot!}\to\mathcal C_{\bot!}\)
and \(Q:\mathcal C_{\bot!}\times\mathcal E_{\bot!}\to\mathcal E_{\bot!}\),
put \(H(B)=\operatorname{FIX}(P(-,B))\),
\(K(B)=Q(H(B),B)\), \(B_* =\operatorname{FIX}(K)\),
\(A_*=H(B_*)\). Their two structure maps make \((A_*,B_*)\)
an initial algebra for \((A,B)\mapsto(P(A,B),Q(A,B))\).
Hence the nested and simultaneous canonical solutions are
isomorphic by actual structure-preserving comparison maps.
The value assertion needs no strictness; the domain initiality
assertion uses the displayed strict regime.
\cite[pp.~73--74; Exercise~5.4.11(4), p.~83]
{AbramskyJungCorrectedExpandedDomainTheory}.
\end{hbtheorem}

\begin{hbtheorem}{thm:p10-domain-invariance-resumptions}{proved=annex:p10.domain-invariance.resumptions}{The source resumption functor and observable recursive behavior}
For a pointed state dcpo \(S\), define
\[
 T(R)=[S\longrightarrow(S\oplus(S\times R))],\qquad
 T(u)(r)=(1_S\oplus(1_S\times u))\circ r
\]
for \emph{strict} continuous \(u\). This is an actual locally
continuous functor on \(\mathsf{DCPO}_{\bot!}\). Its canonical
solution has inverse fold/unfold, initial strict-algebra and
final strict-coalgebra properties. No unrestricted-map extension
is inferred from the coalesced-sum notation.
For a nonbottom state \(s\), define
\[
 \operatorname{return}_s=\operatorname{fold}(x\mapsto\operatorname{inl}(s)),
 \quad
 \operatorname{step}(s,r)=\operatorname{fold}
              (x\mapsto\operatorname{inr}(s,r)).
\]
The continuous functional \(r\mapsto\operatorname{step}(s,r)\)
has a least fixed point whose \(n\)-th approximation exposes
\(n\) successive intermediary states and then an unknown remainder.
Its limit never has a final-return observation; return and step
remain distinct coproduct-branch observations. This example
retains the source's final-state versus state-plus-resumption
interpretation without asserting an unproved programming-language
adequacy theorem \cite[Section~5.1.2, pp.~68--69]
{AbramskyJungCorrectedExpandedDomainTheory}.
\end{hbtheorem}

\subsection{Admissible relations, induction and coinduction}

\begin{hbdefinition}{def:p10-domain-invariance-relations}{Admissibility and typed relation actions}
For finite arity \(n\), \(R\subseteq D^n\) is admissible when
it contains the constantly-bottom tuple and is closed under
coordinatewise suprema of increasing \(\omega\)-chains. Let
\(\mathcal R_n(D)\) be these relations ordered by inclusion.
For strict \(u:D\to E\), write
\(u^{\mathrm{rel}}R=\{(u(x_1),\ldots,u(x_n)):(x_i)_i\in R\}\).
An admissible action for mixed \(F\) assigns
\(F^{\mathrm{rel}}(R,S)\in\mathcal R_n(F(D,E))\).
It must satisfy the typed condition: if
\(u:D_2\to D_1,v:E_1\to E_2\) are strict and
\(u^{\mathrm{rel}}R_2\subseteq R_1\),
\(v^{\mathrm{rel}}S_1\subseteq S_2\), then
\[
 F(u,v)^{\mathrm{rel}}F^{\mathrm{rel}}(R_1,S_1)
                       \subseteq F^{\mathrm{rel}}(R_2,S_2).
\]
It is therefore antitone in the negative relation and monotone
in the positive relation. These actions are not obtained merely
by treating the relations as domain objects and applying \(F\).
\cite[Definitions~5.4.1,5.4.3; Proposition~5.4.4, pp.~80--81]
{AbramskyJungCorrectedExpandedDomainTheory}.
\end{hbdefinition}

\begin{hbtheorem}{thm:p10-domain-invariance-relational}{proved=annex:p10.domain-invariance.relational}{The general admissible-relation comparison principle}
The predicate \(\{u:u^{\mathrm{rel}}R\subseteq S\}\) is
admissible on the strict function dcpo whenever \(R,S\) are
admissible. If \(D\) is the canonical mixed solution and
\[
 \operatorname{unfold}^{\mathrm{rel}}R\subseteq F^{\mathrm{rel}}(S,R),
 \qquad
 \operatorname{fold}^{\mathrm{rel}}F^{\mathrm{rel}}(R,S)\subseteq S,
\]
then \(R\subseteq S\). This follows from the least-identity
endomorphism and an actual invariant admissible predicate.
\cite[Proposition~5.4.2; Theorem~5.4.5, pp.~80--81]
{AbramskyJungCorrectedExpandedDomainTheory}.
\end{hbtheorem}

\begin{hbtheorem}{thm:p10-domain-invariance-actions}{proved=annex:p10.domain-invariance.actions}{Actual constructor actions and variance-compatible composition}
For unary predicates the logical actions on lifting, product,
sum and function space are
\[
 R_{\bot}=\{\bot\}\cup\operatorname{up}(R),\quad
 R\times S,\quad\operatorname{inl}(R)\cup\operatorname{inr}(S),
 \quad [R\longrightarrow S]=\{f:f(R)\subseteq S\}.
\]
Smash and strict function space use the corresponding bottom
quotient and strict-map restriction. For binary relations, lifting
uses pairs whose first coordinate is new bottom or whose two
coordinates are old lifts related by \(R\); products act
coordinatewise; sums use a bottom first coordinate or a pair
in the same branch related by its component relation. Function
space uses the pointwise action
\[
 [R\longrightarrow S]=\{(f,g):(\forall x)\ (f(x),g(x))\in S\},
\]
independent of \(R\). These are admissible actions with their
declared constructor variances. The unary mixed actions are
logical: \(F^{\mathrm{rel}}(R,E)=F(D,E)\); the binary ones
are extensional: \(F^{\mathrm{rel}}(R,\sqsubseteq_E)
=\sqsubseteq_{F(D,E)}\). For
\(H(X,Y)=F(G_1(Y,X),G_2(X,Y))\), their composition is
\[
 H^{\mathrm{rel}}(R,S)=
 F^{\mathrm{rel}}(G_1^{\mathrm{rel}}(S,R),G_2^{\mathrm{rel}}(R,S)).
\]
The binary lifting relation lives on \((D_{\bot})^2\), with
explicit old-element injections and a separate new-bottom clause.
\cite[pp.~81--83]{AbramskyJungCorrectedExpandedDomainTheory}.
\end{hbtheorem}

\begin{hbtheorem}{thm:p10-domain-invariance-induction-coinduction}{proved=annex:p10.domain-invariance.induction-coinduction}{Structural logical induction and extensional simulation}
For a logical unary action, an admissible \(S\subseteq D\)
with \(\operatorname{fold}(F^{\mathrm{rel}}(D,S))\subseteq S\)
equals \(D\). For an extensional binary action, an admissible
\(R\subseteq D^2\) with
\[
 \operatorname{unfold}^{\mathrm{rel}}R
                       \subseteq F^{\mathrm{rel}}(\sqsubseteq_D,R)
\]
is contained in \(\sqsubseteq_D\). Such an \(R\) is a
simulation; \(x\sqsubseteq y\) iff some simulation relates
them. These conclusions use the declared action, not a slogan
that every arbitrary relation is inductive or coinductive.
If a canonical solution is \(\omega\)-algebraic and its
compact basis is given by finite constructor expressions, induction
on those expressions plus admissibility
also establishes a predicate on the whole domain. Without
countable approximation, containing all compact points and
being only omega-chain closed is insufficient.
\cite[Sections~5.4.1--4; Theorems~5.4.7,5.4.9;
Corollary~5.4.10; Proposition~2.2.13, pp.~22,79--83]
{AbramskyJungCorrectedExpandedDomainTheory}.
\end{hbtheorem}

\begin{hbexample}{def:p10-domain-invariance-stream-simulation}{Two recursive descriptions of the same infinite stream}
In the canonical prefix-word domain of
Theorem~\ref{thm:p10-domain-recursive-examples}, fix a letter \(a\).
If \(r,q\) both satisfy
\(\operatorname{unfold}(r)=\operatorname{up}(a,r)\) and
\(\operatorname{unfold}(q)=\operatorname{up}(a,q)\), then
\(\{(\bot,\bot),(r,q)\}\) is an admissible simulation for
the lifted product constructor. Coinduction gives \(r\sqsubseteq q\);
the reversed relation gives \(q\sqsubseteq r\), hence equality.
Annex~\ref{annex:p10.domain-invariance.induction-coinduction}
checks its actual action and admissibility.
\end{hbexample}

\begin{hbtheorem}{thm:p10-domain-invariance-failures}{proved=annex:p10.domain-invariance.failures}{Universal-property and induction boundaries that have counterexamples}
For the identity functor, the two-point pointed chain with
identity structure solves the domain equation but is not initial;
identity and the constant-bottom strict map are distinct
endomorphisms. Its canonical solution is the one-point domain.
Allowing non-strict algebras/maps destroys strict initiality:
every constant map from the one-point domain to a two-point
identity algebra is a homomorphism. The same example with
\(F(X,Y)=Y\) defeats uniqueness of the free mixed comparison
when its backward map is not required strict.
For uncountable \(U\), the predicate of countable subsets of
\(\mathcal P(U)\) contains bottom and every compact element
and is omega-chain closed, but is not the whole domain.
Finally initial algebra and final coalgebra for \(X\mapsto1+X\)
in sets have carriers \(\mathbb N\) and
\(\mathbb N\cup\{\infty\}\); their algebra structures are
not isomorphic. Domain algebraic compactness is not inferred
for an arbitrary category from a recursive equation.
\cite[Exercises~5.4.11(5)--(7), p.~83; pp.~13,18--23]
{AbramskyJungCorrectedExpandedDomainTheory}.
\end{hbtheorem}
